\section{Experiments}
\label{sec:experiments}
\subsection{Shared Six-Dataset Evaluation}
The shared S evaluation fixes the same model, input derivation contract, controls, and eligible test records within each dataset. Table~\ref{tab:shared_results} reports all six datasets together. The 42 common cells are backed by individual predictions; two extra raw-feature cells and ten five-family cells provide incomplete additional axes. One archived mixed-task CARE-PD diagnostic is excluded from the common task. No unexecuted cell is interpreted as a failure.

\begin{table*}[t]
\centering\small
\caption{\textbf{Same-model evidence use and adaptation across all six datasets (S).} Values are macro-F1, recomputed from released predictions with a fixed task label set. LR/SVM values are the committed task-matched supervised references. $N$ is clips/participants. Best LLM view is bold; bold does not imply superiority over references.}
\label{tab:shared_results}
\begin{tabular*}{\textwidth}{@{\extracolsep{\fill}}lrrrrrrr@{}}\toprule
Dataset & $N$ & Chance & ZS summary & ZS z-features & Summary SFT & LR & SVM\\\midrule
ASD-Kinect & 20/20 & .490 & .310 & \textbf{.436} & .310 & .697 & .744\\
ATAXIA & 8/4 & .488 & .273 & \textbf{.333} & \textbf{.333} & .873 & .873\\
CARE-PD & 665/24 & .250 & .321 & .133 & \textbf{.445} & .659 & .650\\
GMFCS & 255/149 & .250 & .265 & .150 & \textbf{.467} & .579 & .571\\
PARKINSONISM & 29/5 & .494 & .396 & \textbf{.757} & .542 & .561 & .574\\
PATHOGAIT & 1440/2 & .166 & .069 & .139 & \textbf{.451} & .881 & .897\\\bottomrule
\end{tabular*}
\end{table*}

\paragraph{Zero-shot evidence use.}
Summary zero-shot macro-F1 is near stratified chance or below it on several datasets, but performance alone does not establish whether a model used the input. CARE-PD exceeds shuffled evidence by 0.065 macro-F1 (archived participant-cluster 95\% interval $[0.007,0.118]$, paired bootstrap tail area 0.029). This is evidence of clip-specific responsiveness in that tested setting, while the 0.321 score remains far below its 0.659 LR reference. GMFCS's summary advantage over shuffled evidence is 0.052 (tail area 0.089), and reversing candidate order raises macro-F1 to 0.322. The latter sensitivity cautions against equating plausible accuracy with stable evidence grounding.

The five-family axis supports the accepted finding of near-chance L agreement ($\kappa=-0.069$ to 0.035), but these are protocol-scoped results. The available S CARE-PD cells have $\kappa=0.099$--0.130; ASD-Kinect ranges from $-0.100$ to 0.200. Thus, neither the original nor expanded evidence justifies a universal claim that zero-shot LLMs cannot use pose-derived evidence. Legacy GMFCS cells affected by unresolved split provenance are retained only for audit in the supplement and are excluded from the main six-dataset comparisons.

\subsection{Input Representation and Recoverable Signal}
\label{sec:modality}
Input views can change prediction distributions without consistently restoring agreement. In the accepted L view experiments, $\kappa$ remains between $-0.095$ and 0.057; this range must not be applied to S. Under S, PARKINSONISM's named z-feature view reaches 0.757 macro-F1 and $\kappa=0.520$, versus 0.396 macro-F1 for summaries. Its archived subject-cluster interval is $[0.543,0.906]$, based on only 29 clips from five participants. The point estimate exceeds the matched linear references, but this small test does not establish population-level superiority. CARE-PD shows the opposite representation preference: summary macro-F1 is 0.321 versus 0.133 for z-features.

On ASD-Kinect, summary and z-feature scores are 0.310 and 0.436, both below the matched linear point estimates. ATAXIA's eight clips cannot support a powered model ranking. PATHOGAIT is low under both zero-shot views; GMFCS summaries outperform its z-feature encoding. These six rows jointly show representation dependence rather than a universally best input view.

The same-feature references establish recoverable supervised signal, with LR macro-F1 from 0.561 to 0.881 under S. Reference strength and uncertainty remain dataset-specific: ASD-Kinect LR has a wide interval and an archived permutation value of 0.09. They do not show that supervised and zero-shot learning have equivalent information access. Per-class recall and confusion matrices in the supplement expose concentrated outputs that aggregate accuracy can conceal.

\subsection{Adaptation Recovers Dataset-Specific Mappings}
\label{sec:adaptation}
Summary-SFT changes macro-F1 by 0.000, 0.061, 0.124, 0.202, 0.146, and 0.382 for ASD-Kinect, ATAXIA, CARE-PD, GMFCS, PARKINSONISM, and PATHOGAIT, respectively. GMFCS improves to 0.467; its archived paired participant-bootstrap tail area is 0.001. CARE-PD improves to 0.445, but its corresponding tail area is 0.43. Both remain below matched LR. PARKINSONISM SFT remains below its zero-shot z-feature result. PATHOGAIT improves substantially but remains below its capture-quality-free linear references; two participants prevent a population-level gain claim. ASD-Kinect shows no summary-SFT macro-F1 gain, despite a change in the class receiving most predictions.

The accepted-protocol adaptation ladder is preserved as a separate, explicitly scoped experiment. Few-shot ATAXIA improves from 0.406 to 0.697 macro-F1. Within the same 7B model, ATAXIA summary agreement rises from $\kappa=0.046$ to 0.507, and feature-SFT reaches 0.762 macro-F1 versus LR 0.763. Its participant comparison does not detect a difference ($p=0.97$); this is not an equivalence test. PATHOGAIT feature-SFT reaches 0.687 versus LR 0.662, but the direction reverses for one of its two participants. These values retain the accepted in-distribution recovery finding; they are not S reruns or evidence of general diagnostic knowledge.

One L PATHOGAIT feature, valid-frame fraction, carries capture-related information. Removing it reduces LR/SVM macro-F1 from 0.662/0.644 to 0.564/0.560. The original feature-SFT was not rerun with that feature removed. The common S evaluation excludes it for all six datasets, but also changes splits and feature construction, so it cannot isolate that single feature's causal contribution. We therefore retain the L adaptation result only as learning from its supplied feature set, and make no kinematics-only superiority claim.

\subsection{Transfer and Qualitative Diagnostics}
\label{sec:transfer}
In the tested L ATAXIA--PATHOGAIT transfer probe, in-domain $\kappa$ is 0.549 and 0.917 for the two directions, while target-domain $\kappa$ is 0.057 and 0.000. PATHOGAIT is binarized to an approximately 11\%/89\% target distribution; the resulting binary accuracy cannot be compared with six-class scores. This supports only the accepted conclusion that reliable transfer was not observed on this harmonized pair. Six-dataset membership does not imply task-compatible transfer pairs.

The supplement adds prediction-backed success/failure cases for each of the six datasets and separates them from source-publication images in Fig.~\ref{fig:method_pipeline}. Those images illustrate resource appearance, not model evidence. Full S pose/input records are absent from the upstream snapshot, so a verified, clip-matched pose--input--output panel cannot yet be reconstructed. We do not substitute an unrelated frame or invented feature description. The release records the missing assets needed to complete that reviewer request.

\subsection{Release and Limitations}
\label{sec:release}
The independent project repository provides six-dataset prediction logs, scoring/plotting code, the actual runner and SFT configuration, a table-to-file map, checksums, and source-access/license notes. It regenerates the shared LLM point metrics and figures from committed outputs. Source poses, complete train/validation/test input layers, adapters, and training logs are not included; full training/inference reproduction requires these assets from their authorized source. Dataset licenses are not replaced by a blanket repository license. Original RGB recordings are not redistributed, and GMFCS public access is to de-identified trajectories rather than clinical video.

Small test populations, performed-gait provenance, source simulation uncertainty, unequal class support, and unresolved training-artifact availability limit interpretation. The subject-disjoint S contract addresses the \emph{design} of the old GMFCS overlap concern, but its full split inputs must still be recovered for independent verification; no newly executed leakage-free rerun is claimed. These boundaries are disclosed with the reproducible result checks instead of treating missing evidence as a successful experiment.
